\section{はじめに}
近年情報通信の技術は発展しており,社会を生きるうえで情報のやり取りを行うのは必須となっている.
その中でも無線通信の需要が高まっており,スマートフォンという身近なものだけでなく医療機器や教育用端末などの多様な場面でも使用されている.
総務省の調査によると2022年時点で自宅無線LANを利用している人は84.9 \%にも及んでいる\cite{soumusyou}.
このような情報通信の重要性は学校の教育にも浸透しており小中学校では情報の授業が必修になっており情報通信を学ぶ機会が設けられている\cite{kyouiku}.
今まで以上に情報への深い興味と理解が必要になってきている.
一方で今主流である無線通信の多くが赤外線通信や音響通信のように人には可視・可聴できないものである.
赤外線通信は携帯電話の連絡先交換や写真のデータ交換をするときなどに使用されているが,赤外線のため不可視であり通信が成功したかどうかを直感的に判断しにくい\cite{sekigaisen}.
また音響通信は人には聞こえない高域帯用いており,人間には聞こえない\cite{onkoyu}.
モールス信号は人に聞こえるが,聞き取るには練習が必要である.
そのため通信過程の視認性や教育的わかりやすさにかけるという課題がある.
一方で視覚的に動作がわかりやすい教材としてNHK放送のピタゴラスイッチという番組がある.
この番組では球が坂を転がったり,ドミノが倒れて装置を動かしたりなどの物理的動作で「ピ」という文字を伝達している\cite{pitagora}.
このような装置は可視で通信過程を視認でき面白いが,あくまでも遊びや仕組みの面白さを重視しており情報伝達手段としては生かされていない.


そこで本提案では赤外線通信と音響通信のように情報を伝達できる機能とピタゴラスイッチのように視覚的に情報伝達している様子がわかりやすい動作を組み合わせた,情報伝達の様子が直感的にもわかりやすい非電波の無線通信の情報伝達装置を提案する.
提案する装置は以下の3種類の通信方式からなっている.
すべてASCIIコードを使用しており2,4,8進数で情報のやり取りを行っている.
1つ目はスピーカとマイクを使用した音のシステムである.
この装置は可聴範囲内の周波数の音をスピーカーから送信し,マイクで受け取ることで文字を構成する.
2つ目はLEDと照度センサを使用した光のシステムである.
この装置は可視光線をLEDで点滅させ明暗のパターンによって文字を構成する.
3つ目は球と照度センサを使用した球のシステムである.
この装置は球をモーターで分岐させ照度センサの上に通し明暗を作っている.
これはピタゴラスイッチ的な球の動きと通信を組み合わせた新しいシステムである.


本提案ではこれらの3つの情報伝達装置について1文字あたりの伝達速度と文字認識の伝達精度・書き取り精度という3つの指標で評価を行う.
その結果,評価目標は音の伝達精度,伝達速度,光の書き取り精度,球の伝達速度以外は目標を上回る事ができた.
ここから待機時間の調整などを行えば教育に活かせることがわかった.
